\documentclass{article}[12pt]
\usepackage{german}

\title{Simulation hamiltonscher Dynamiken auf wechselwirkenden
Quantensystemen} 

\author{Pawel Wocjan\\Institut f\"ur Algorithmen und Kognitive Systeme\\
Universit\"at Karlsruhe}

\date{12. Juli 2002}

\begin{document}
\maketitle 

Eine Wechselwirkung zwischen $n$ Quantensystemen l"asst sich dazu
verwenden, Zeitentwicklungen zu simulieren, die v"ollig anderen
Hamiltonoperatoren entsprechen. Dabei wird die nat"urliche
Zeitentwicklung durch Sequenzen von lokalen Operationen auf den
Komponenten unterbrochen. Die wechselseitige Simulierbarkeit
verschiedener Hamiltonoperatoren erlaubt den Vergleich der
"`Berechnungsm"achtigkeit"' von Quantencomputern, die unterschiedliche
physikalische Wechselwirkungen ausnutzen.

Der Vortrag hat die Komplexit"at und den Entwurf effizienter
Simulationsalgorithmen zum Thema. Die Bestimmung des minimalen Zeit-
und Schrittaufwand l"auft auf ein konvexes Optimierungsproblem
hinaus. Obere und untere Schranken f"ur den Simulationsaufwand lassen
sich gewinnen, indem man die Hamiltonoperatoren durch Graphen
beschreibt und die Spektren der zugeh"origen Adjazenzmatrizen
vergleicht. Effiziente Simulationsalgorithmen basieren auf Methoden
aus der Darstellungstheorie endlicher Gruppen, der Designtheorie und
der konvexen Optimierungs\-theorie. F"ur ein Paar interessante
Simulationsalgorithmen konnte mit den hergeleiteten unteren Schranken
ihre Optimalit"at bewiesen werden.


\end{document}
\documentclass{article}[12pt]
\usepackage{german}

\title{Simulation hamiltonscher Dynamiken auf wechselwirkenden
Quantensystemen} 

\author{Pawel Wocjan\\Institut f\"ur Algorithmen und Kognitive Systeme\\
Universit\"at Karlsruhe}

\date{12. Juli 2002}

\begin{document}
\maketitle 

Eine Wechselwirkung zwischen $n$ Quantensystemen l"asst sich dazu
verwenden, Zeitentwicklungen zu simulieren, die v"ollig anderen
Hamiltonoperatoren entsprechen. Dabei wird die nat"urliche
Zeitentwicklung durch Sequenzen von lokalen Operationen auf den
Komponenten unterbrochen. Die wechselseitige Simulierbarkeit
verschiedener Hamiltonoperatoren erlaubt den Vergleich der
"`Berechnungsm"achtigkeit"' von Quantencomputern, die unterschiedliche
physikalische Wechselwirkungen ausnutzen.

Der Vortrag hat die Komplexit"at und den Entwurf effizienter
Simulationsalgorithmen zum Thema. Die Bestimmung des minimalen Zeit-
und Schrittaufwand l"auft auf ein konvexes Optimierungsproblem
hinaus. Obere und untere Schranken f"ur den Simulationsaufwand lassen
sich gewinnen, indem man die Hamiltonoperatoren durch Graphen
beschreibt und die Spektren der zugeh"origen Adjazenzmatrizen
vergleicht. Effiziente Simulationsalgorithmen basieren auf Methoden
aus der Darstellungstheorie endlicher Gruppen, der Designtheorie und
der konvexen Optimierungs\-theorie. F"ur ein Paar interessante
Simulationsalgorithmen konnte mit den hergeleiteten unteren Schranken
ihre Optimalit"at bewiesen werden.


\end{document}
